\paragraph{Finite evaluations for the geometric index.}
\label{par:slice-evaluation-net}

A colength in the product module can be estimated by restricting sections to an evaluation set. The norm used for that restriction is fixed on the same run.

Work on the run at weight \(\kappa\) with coefficient field
\(K_0\), and write \(V_\kappa\) for its rational cusp space.
The analytic evaluations below use the uniformization of this
space on the level of the run. In particular they can be
restricted to any of the cusp subspaces tested by the index.
Put \(Z_b=frL\), where \(2f=[E:\Q]\) for the CM field \(E\), \(r\)
is the rank and \(L=\log b\).

\begin{proposition}[Finite evaluations on the selected run]\label{prop:slice-evaluation-net}
Fix a run and lift in the form coordinates of
Definition~\ref{def:common-row-lattices}. Use its rational cusp space
\(V_\kappa\) of weight \(\kappa\) over \(K_0\), with \(Z_b=fr\log b\)
and the
archimedean and level sizes of \eqref{eq:slice-net-run-sizes}.
For every embedding \(\sigma:K_0\hookrightarrow\C\) there is a nonempty finite set
\(\mathcal N_\sigma\) of scalar evaluation functionals, linear over \(\sigma\).
If \(m_\sigma=|\mathcal N_\sigma|\), they satisfy
\begin{equation}\label{eq:slice-net-size}
 \frac14\|F\|_{\sup,\sigma}\le
    \max_{y\in\mathcal N_\sigma}|e_{\sigma,y}(F)|
      \le\|F\|_{\sup,\sigma},\qquad
 \log(16m_\sigma)\le1041Z_b^3 .
\end{equation}
Complex conjugation identifies the sets of functionals at conjugate embeddings, and the simultaneous evaluation map is injective after extending
scalars from \(K_0\) to \(\C\) at each embedding.

\end{proposition}

\begin{proof}
\proofstep{Reduction to a self-dual Hermitian lattice.}
First use the embeddings sending the marked imaginary
quadratic root to the chosen root. The components there are
framed Hermitian domains with the character, principal
congruence subgroup and scalar weight of the run. In each
component write \(\Lambda\) for the symplectically self-dual
period lattice with pairing
\(\operatorname{Tr}_{E/\Q}(\xi H)\), for \(H\) its Hermitian
form and \(\bar\xi=-\xi\). Divide \(\Lambda\) by the fixed
rank-one fractional ideal \(I\) of the principal datum:
\(L'=I^{-1}\Lambda\). The rank-one trace pairing is perfect,
in the sense that
\[
 \operatorname{Tr}_{E/\Q}(\xi a\bar b x)\in\mathbb Z
 \quad(a,b\in I)\qquad\Longleftrightarrow\qquad x\in\mathcal O_E.
\]
For every \(u\in E^r\), test the trace pairing of \(Iu\) with
\(I L'=\Lambda\). The rank-one identity gives
\[
 \{u:H(u,L')\subseteq\mathcal O_E\}
   =\{u:\operatorname{Tr}_{E/\Q}(\xi H(Iu,\Lambda))
                        \subseteq\mathbb Z\}
   = I^{-1}\Lambda = L'.
\]
The second equality uses the self-duality of \(\Lambda\).

Take positive integers
\(N_A,N_C,m_0\) with \(N_AI\subset\mathcal O_E\),
\(N_C\in I\), and \(2m_0\xi\in\mathcal O_E\).
Thus \(N_A\Lambda\subset L'\) and \(N_C L'\subset\Lambda\).
The principal level \(g\) is the product of the two prescribed levels. Write \(T_v(\mu)=H(\mu,v)v\) in the convention
linear in \(\mu\). For \(a\in\mathcal O_{E^+}\) the nilpotent
part of the transvection with parameter \(Ma\),
\(M=gN_AN_Cm_0\), takes values in \(g\Lambda\) because
\[
  N_A H(\Lambda,v)\subseteq\mathcal O_E,\qquad
  N_C v\in\Lambda,\qquad 2m_0\xi\in\mathcal O_E .
\]
Since \(v\) is isotropic, \(T_v^2=0\), so the inverse has the
negative of the same nilpotent part. The transvection preserves
\(\Lambda\) and is trivial modulo \(g\). It belongs to the
principal subgroup and the section character is trivial
there. The same construction applies with the opposite
choice of \(\xi\) on a component.

\proofstep{The size conditions on the selected run.}
Here are the bounds used on the selected run. Write \(d_E\)
for the absolute discriminant, \(k_0=2|\tau(\xi)|\) on the
chosen half, \(\lambda_{\tau,h}\) for the frame weight,
\(S=\sum_{\tau,h}|\lambda_{\tau,h}|\), and choose
\(W>0\) at least the sum of the spreads
\(\max_h\lambda_{\tau,h}-\min_h\lambda_{\tau,h}\).
The bounds on the selected parameters give
\begin{equation}\label{eq:slice-net-run-sizes}
 \log N_A\le4L,\quad \log N_C\le2L,\quad
 \log m_0\le4L,\quad \log g\le19L,\quad 0<k_0\le2 .
\end{equation}
Here \(N_C\) is the norm of the ideal used in the rank-one
datum. Thus \(\log M\le30L\).

At \(\kappa\) we have
\(f\le r\le n_\kappa\), \(f\ge16\),
\(\log f\le2L\), \(\log n_\kappa\le L\) and
\(\log d_E\le4fL\). The differential and central parts of the
weight give \(W=D_F+fn_\kappa+1\) and
\(S\le D_F+frn_\kappa/2\), where \(D_F\) is the sum of the
nonnegative flag degrees. The size inequalities for this weight
give \(\log(1+D_F)\le2L\). Since \(f\le r\le n_\kappa\le e^L\) and \(1+D_F\le e^{2L}\),
\[
 W\le2e^{2L}\le e^{3L},\qquad
 1+S\le e^{2L}+\tfrac12 e^{3L}\le e^{3L}
\]
in the selected range \(L\ge2\).
These bounds use the fixed run's prescribed central weight,
level and field. The two run orientations satisfy them
separately. \label{cor:slice-evaluation-net-1}

\proofstep{Bounding the section away from short isotropic directions.}
For a cusp section on one component, write \(P(x)\) for its
scalar in the frame \(x\), \(n(x)\) for the product of the
absolute flag-minor ratios to powers
\(-\lambda_{\tau,h}/2\), and \(U=\sup_x|P(x)|n(x)\).
This is a finite supremum by the cusp condition. Let \(Q_{x,\tau}\)
be the positive majorant of \(H_\tau\) obtained by reversing
the negative part in the orthogonal flag of \(x\). Set
\[
 \varepsilon_* = \min\left(1,
       \frac{\pi}{M(fW+2)k_0\log 2}\right).
\]
If \(U>0\), take a frame \(x\) with \(|P(x)|n(x)>U/2\).
Every nonzero \(v\in L'\) then has
\(Q_{x,\tau}(v)\ge\varepsilon_*\) at some \(\tau\).

To prove this, assume all the majorants are smaller.
The signed form is dominated by the positive majorant, so
\(|\tau H(v,v)|\le Q_{x,\tau}(v)<\varepsilon_*\le1\) for every
chosen embedding. Because \(H(v,v)\in\mathcal O_{E^+}\),
\[
      |\mathrm N_{E^+/\Q}H(v,v)|
          =\prod_\tau |\tau H(v,v)|<1 .
\]
This rational integer is zero. Thus the small vector is isotropic
and the rank-one transvection applies.
On the complex transvection through \(x\) in direction \(v\),
use the same parameter \(s\) at every \(\tau\).
The condition on the \(\tau\)-flag is
\(1+k_0Q_{x,\tau}(v)\Im s>0\).

On the line \(\Im s=-t\), the upper bound \(Q_{x,\tau}(v)\le\varepsilon_*\)
makes the margin explicit:
\[
 1-k_0t Q_{x,\tau}(v)\ \ge\ 1-k_0t\varepsilon_*=\tfrac12,
 \qquad t=(2k_0\varepsilon_*)^{-1}.
\]
This criterion and the norm comparison follow by taking the
leading minors of the rank-one change of the frame Gram
matrix. On \(\Im s=-t\), with \(t=1/(2k_0\varepsilon_*)\),
the ratio of the original metric factor to the new one is
at most \(2^{W/2}\), by telescoping the minor ratios with
the weight differences. The whole function \(P(x(s))\) is
holomorphic for \(\Im s>-2t\), and is \(M\)-periodic there.

\proofstep{The cusp expansion and the maximum principle.}
To see that it vanishes at the upper end, use the cusp
expansion in the integral isotropic direction \(N_Cv\).
The transvection is quadratic in its isotropic direction. Thus its parameter \(s'\) for \(N_Cv\) has to satisfy
\[
  N_C^2s'=s,\qquad s'=s/N_C^2.
\]
The Fourier frequencies are nonnegative in each real embedding
by boundedness in the normalized metric: the metric changes
at most polynomially in the positive imaginary directions,
so a negative frequency would grow exponentially there.
The zero frequency vanishes by cuspidality. The expansion
is absolutely convergent at the original frame. Along the
positive diagonal, each remaining exponential tends to zero
and is bounded by one in modulus. This proves the upper-end
vanishing. Put 
\[
  z=e^{2\pi i s/M}
\]
 and \(R=e^{2\pi t/M}\). Periodicity gives
a holomorphic function \(F(z)=P(x(s))\); the upper-end limit gives
\(F(0)=0\).

Apply the maximum principle to \(F(z)/z\) on \(|z|\le R\):
\[
 |F(1)|\le R^{-1}\max_{|z|=R}|F(z)|
        =e^{-2\pi t/M}\max_{\Im s=-t}|P(x(s))|.
\] The result is
\[
 |P(x)|n(x)\le e^{-2\pi t/M} U\,2^{W/2}.
\]
Since \(2\pi t/M\ge(fW+2)\log 2\), the right side is strictly
less than \(U/2\). Thus the set of frames whose squared first
minimum for \(\sum_\tau Q_{x,\tau}\) on \(L'\) is at least
\(\varepsilon_*\) contains a point at which each given nonzero component section nearly attains its supremum. This threshold is independent of
the section. \label{cor:slice-evaluation-net-2}

\proofstep{A finite grid of reduced frames.}
We now construct one finite set on this part with larger weights. Use
\[
 B=\max\left(1,\frac{(2fr)^4d_E^2}{4\varepsilon_*}\right),
 \qquad
 \eta=\frac{1}{1000r^2(1+S)}.
\]
The Euclidean metric \(\sum_\tau Q_{x,\tau}\) has real
dimension \(2fr\). Its dual lattice is obtained by applying
twice the isometric sign involution of the majorant to the
trace dual of \(L'\). Multiplication by \(d_E\) sends that
trace dual into \(L'\), by Hermitian self-duality and the
different. The Euclidean dual first minimum is therefore
at least \(2\sqrt{\varepsilon_*}/d_E\). The real transference bound gives
\[
 \lambda_{2fr}(L')
 \le\frac{2fr}{\lambda_1((L')^*)}
 \le\frac{(2fr)d_E}{2\sqrt{\varepsilon_*}} .
\]

The vectors attaining the real rank span the \(E\)-space, since
any proper \(E\)-subspace has smaller real dimension. Select \(r\)
independent vectors over \(E\); the square of the displayed bound
is at most \(B\). The larger factor \((2fr)^4\) also suffices in the transference estimate. \cite{banaszczyk}
In an orthonormal flag their Hermitian Gram entries
are bounded in absolute value by \(B\) at all embeddings.
Orthonormalization here is right translation by an upper
triangular frame change and preserves the normalized value of the section.

There are at most \((1+2B)^{2fr^2}\) such Gram matrices.
For each entry, Minkowski space packing in the maximum of the
complex moduli gives \((1+2B)^{2f}\), since distinct integers
of \(E\) differ by an element with some conjugate of modulus
at least one.
Using the entry bound independently at the \(r^2\) matrix positions gives
\[
  \big((1+2B)^{2f}\big)^{r^2}=(1+2B)^{2fr^2}.
\]
 Fix a Gram matrix \(A\).

To classify the
short bases modulo the principal congruence group, pair
\(N_A\Lambda\) with each basis, and read the pairing vectors
modulo the image of \(gN_AN_CA\) on \(\mathcal O_E^r\).
The values on \(r\) generators of the corresponding finite
quotient of \(\Lambda\) suffice. Such generators exist because
a finite quotient by a nonzero scalar of a projective module
of rank \(r\) is free of rank \(r\) over the quotient ring.
With \(C=gN_AN_CA\), the required scalar is \(\det C\). Indeed
\[
 (\det C)\,\mathcal O_E^r
 =C\,\operatorname{adj}(C)\mathcal O_E^r\subseteq C\mathcal O_E^r .
\]
If the residues agree, the isometry taking one basis to
the other sends every \(x\in\Lambda\) to a vector differing
from \(x\) by an element of \(gN_C\operatorname{span}_{\mathcal O_E}v\).
Apply the same argument to the inverse to obtain equality of
lattices. The isometry belongs to the principal subgroup.
The size of the pairing quotient is

\[
  |N_{E/\Q}\det(gN_AN_CA)|\le(rgN_AN_CB)^{2fr},
\]
by the determinant expansion. The number of basis types per
Gram matrix is at most its \(r\)-th power.

For each representative basis \(v\), consider all orthonormal
frames \(x\) in which its columns have squared majorant length
at most \(B\). At \(\tau\) put \(C_\tau(x)=x_\tau^{-1}v_\tau\).
Its entries have modulus at most \(\sqrt B\), and its
determinant has fixed modulus
\(d_\tau=|\tau(\det A)|^{1/2}\).
Since \(\prod_\tau d_\tau\ge1\), one can use mesh sizes
\[
 \frac{\eta d_\tau}{2r^{r+1}(\sqrt B)^r}
\]
on the real and imaginary coordinates at \(\tau\).

On an occupied cell the difference of the two coefficient
matrices is bounded entrywise by the mesh, up to the factor
\(\sqrt2\) from its real and imaginary coordinates. The identity
\[
 x_\tau^{-1}y_\tau-1
   =(C_\tau(x)-C_\tau(y))\,
        \frac{\operatorname{adj}C_\tau(y)}{\det C_\tau(y)}
\]
and the bound \((r-1)!(\sqrt B)^{r-1}\) on each cofactor
give an entrywise error at most \(\eta\). Choose one
actual frame per occupied cell. Write \(C=6r^2(rB)^r/\eta\). There are \(2r^2\) real coordinates at
each of the \(f\) embeddings, each with at most \(C/d_\tau\) cells.
Their product is bounded by
\[
 \prod_\tau (C/d_\tau)^{2r^2}
 = C^{2fr^2}\bigl(\prod_\tau d_\tau\bigr)^{-2r^2}
 \le C^{2fr^2}.
\]

\proofstep{Stability of section values on the grid.}
For an orthonormal frame, an entrywise perturbation of radius
\(10\eta\) stays in the domain. Indeed the principal
Gram minors differ from the signed unit minors by perturbations
of the identity. With entry radius \(u\le10\eta\), the Gram
perturbation has entries bounded by \(2u+ru^2\le3u\).
After multiplication by the diagonal signs on each principal
block its operator norm is at most \(3ru\). The singular
values lie between \(1-3ru\) and \(1+3ru\), so the log
of the absolute determinant is bounded by \(9r^2u\). The signed determinant is positive by continuity along the
segment from the orthonormal frame. If \(D_{\tau,h}\) denotes
the absolute leading minor and \(D_{\tau,0}=1\), the metric
has logarithm
\[
 \log n=-\tfrac12\sum_{\tau,h}\lambda_{\tau,h}
             (\log D_{\tau,h}-\log D_{\tau,h-1}).
\]

Applying the minor estimates to both terms gives
\(|\log n|\le9r^2uS\le0.09\). In particular \(n>0.91\)
on the ball, so \(|P|\le U/0.91\) there.
The Schwarz bound on each complex line through the origin,
applied to the difference from the value at the orthonormal
frame, changes the scalar by at most \(0.21U\) at radius
\(\eta\). Thus a representative of a frame of norm greater
than \(U/2\) has normalized norm at least
\(0.91(1/2-0.21)U>U/4\).
Both isometry and orthonormal flag change preserve the norm
before this perturbation. A component with an empty set of indices with larger weights can use one arbitrary frame, since then every cusp
section on it vanishes. \label{cor:slice-evaluation-net-3}

\proofstep{Counting the grid and the components.}
The resulting number \(N\) on a component satisfies
\[
 \log N\le
   2fr^2\log(3N_AN_Cgr B^2)
   +2fr^2\log(6r^2/\eta)+2fr^3\log(rB).
\]
The run bounds give \[
 -\log\varepsilon_*\le41L,\qquad \log B\le65fL,
\] and bound the three logarithms
on the right by \(223fL,20fL,66fL\), respectively.
After adding \(\log16\) the estimate is at most \(1000Z_b^3\).

For the number of components, first pass from the period
lattices to Hermitian self-dual lattices as above. They all
lie in the same rational Hermitian space. In fact, if the
determinant lattice in a chosen basis has fractional ideal
\(J\), self-duality gives 
\[
  (\det H)J\bar J=\mathcal O_E.
\]
The ratio of two Hermitian determinants is therefore a local
norm at every inert finite place, where its valuation is even
and unit norms are surjective.
At an inert unramified prime, for a uniformizer \(\pi\) fixed by conjugation and a unit \(u\), the determinant norm has
\[
  v\big(N(\pi^k u)\big)=2k .
\]
 Every element is a norm at a
split place, and the selected extension \(E/E^+\) is unramified
at every finite place. Equal signatures give equal determinant
signs at the real places. Landherr's classification by signature
and local determinant classes gives the asserted isometry.
\cite{landherr}\label{cor:slice-evaluation-net-4}

An integral self-dual lattice in this space has a majorant
with squared first minimum at least \(f\). To see this, choose
a complete flag in a maximal rational isotropic subspace
and append its annihilators in reverse order. The annihilator
grades pair perfectly with the opposite isotropic grades
on the lattice. Indeed all the intersections defining the
flag are saturated. A functional on one grade extends to
the lattice, vanishing on the preceding flag term, and
self-duality realizes it by a vector in the corresponding
annihilator.
For successive saturated isotropic stages \(I\subset J\), self-duality gives the perfect Hermitian pairing between the two grades
\[
  (I^\perp/J^\perp)\times(J/I)\longrightarrow{\cal O}_E,\qquad ([x],[y])\longmapsto h(x,y).
\]
 In a Witt basis the paired line ideals are
therefore \(J_i\) and \(\bar J_i^{-1}\).

At each chosen
embedding give these lines squared metric factors
\((\mathrm N J_i)^{-1/f}\) and \((\mathrm N J_i)^{1/f}\).
They are reciprocal, as required for a majorant on a
hyperbolic pair. Give the anisotropic middle space an
orthogonal majorant at each embedding. For a nonzero vector \(v\), take its first nonzero
graded image in the increasing flag. If that image is \(a\)
in a line ideal \(J\), the product of its squared lengths
at one embedding from each conjugate pair is
\[
 \prod_\tau\frac{|\tau(a)|^2}{(\mathrm N J)^{1/f}}
      =\frac{|\mathrm N_{E/\Q}a|}{\mathrm N J}\ge1 .
\]
This applies to the opposite grade \(\bar J_i^{-1}\) as well,
since its norm gives the reciprocal metric factor.

If the
image is in the anisotropic middle grade, its Hermitian
square is a nonzero element of \(\mathcal O_{E^+}\); the
product of the absolute values of its conjugates is at
least one. Each absolute value is at most the squared
majorant on this grade. In either case the full squared
majorants dominate these lengths, so
\(\prod_\tau Q_\tau(v)\ge1\) and AM--GM gives
\(\sum_\tau Q_\tau(v)\ge f\).
Transference at this majorant gives bounded bases with the
above value of \(B\) with \(\varepsilon_*\) replaced by \(f\);
call it \(B_0\). A Gram matrix \(A\) and the image of the
lattice under pairing with its basis determine its isometry
class. In the integral pairing coordinates an image \(M\) satisfies
\[
             \operatorname{im} A\subseteq M\subseteq\mathcal O_E^r .
\]
The finite-module argument counts \(M/\operatorname{im}A\)
by at most \(r\) generators.

This generator count gives at most \(((1+2B_0)rB_0)^{2fr^2}\) unmarked lattices. Evaluating the
level marking on \(r\) generators gives at most
\((g^\sharp)^{2fr^2}\) markings per lattice. Thus the number of
components is bounded by
\[
          ((1+2B_0)rB_0 g^\sharp)^{2fr^2}.
\] Here \(g^\sharp\) is the level for the full
marking and \(\log g^\sharp\le30L\).
Polarized complex torus uniformization identifies equal
marked classes. In particular the components number at most
\(\exp(41Z_b^2)\): use
\(\log B_0\le (12+8f)L\), \(1+2B_0\le e^L B_0\) and
\(2(56+16f)\le41fL\) for \(f\ge16, L\ge2\).
Taking the disjoint union over components gives
\[
 \log(16m_\sigma)\le1000Z_b^3+41Z_b^2\le1041Z_b^3.
\]
\label{cor:slice-evaluation-net-5}

\proofstep{Compatibility of scalar evaluation at the test frames.}
At a selected frame \(y\) put
\(e_{\sigma,y}(F)=l_\kappa n(y)P_F(y)\), where
\(l_\kappa>0\) is the fixed Lie-frame factor of the
run and \(P_F\) is the scalar representing \(F\) in these coordinates. The supremum defining
\(\|F\|_{\sup,\sigma}\) uses this same factor. This proves
\eqref{eq:slice-net-size}. These are complex frames. The
field of definition used in the index is \(K_0\) for the
sections, maps and integral lattices. Each evaluation functional is \(\sigma\)-semilinear on this \(K_0\)-space. The height and
local chart construction both use this original field.

\proofstep{Injectivity of the evaluation maps.}
The map to analytic sections is injective after extension to
\(\C\). Let \(A=K_0[t_1,\ldots,t_n]/I\) be a test
algebra, where \(I\) is radical, and let \(y\) be a test
frame over \(A\). For each \(K_0\)-algebra map

\[
  \phi:A\to\C
\]
 over \(\sigma\), base change of the
section rule and complex uniformization identify
\(\phi(F(y))\) with the analytic value at the induced complex frame times a scalar independent of \(F\).
Suppose \(F_1,\ldots,F_h\) are independent over \(K_0\)
and their analytic realizations satisfy \(\sum_i c_iP_{F_i}=0\).
For a fixed test frame \(y\), the tensor
\[
                    b_y=\sum_i c_i\otimes F_i(y)
                    \in\C\otimes_{K_0,\sigma} A
\]
evaluates to zero under every complex point, since the
uniformization scalar at that point is independent of \(i\). This
algebra is reduced because \(K_0\) has characteristic
zero and is perfect. It is of finite type over \(\C\), so
the Nullstellensatz says that the intersection of the
kernels of its complex points is zero. \cite{test-algebra-base-change}
Let \(V\) be the \(K_0\)-span of the coefficients \(c_i\) in
\(\C\), using \(\sigma\) for the scalar action.

The inclusion
\(V\hookrightarrow\C\) remains injective after tensoring with
\(A\) over the field \(K_0\), so the relation is already zero
in \(V\otimes_{K_0}A\). For every dual basis element \(\theta\in V^\vee\),
\[
          0=(\theta\otimes1)b_y=\sum_i\theta(c_i)F_i(y)
                                  \qquad\text{for every }y.
\]
The section \(\sum_i\theta(c_i)F_i\) is therefore zero,
and independence forces all its coefficients to vanish. The dual basis separates
points of \(V\), so all \(c_i\) vanish.
The net controls the supremum for complex linear combinations of these sections as well, since these sums are still holomorphic
cusp sections with the same transformation law. Thus simultaneous evaluation on the finite set is injective over \(\C\). At the other
embeddings use the conjugates of the evaluation functionals on the chosen half.
The two halves use the same supremum by the convention for
the marked imaginary root, and conjugation preserves both
the norm bounds and complex linear independence. \label{cor:slice-evaluation-net-6}

\end{proof}

Define the evaluation norm on the selected weight by
\[
 \|F\|_{2,\sigma}^2
    =\frac1{m_\sigma}\sum_{y\in\mathcal N_\sigma}|e_{\sigma,y}(F)|^2 .
\]
Put \(P_d=1041Z_b^3(1+\log|\kappa|)\).

\begin{proposition}[Evaluation norm and restricted rank]\label{prop:slice-evaluation-norm}
This is a conjugation-compatible Hermitian norm with
\[
 \|F\|_{2,\sigma}\le\|F\|_{\sup,\sigma}
       \le 4\sqrt{m_\sigma}\,\|F\|_{2,\sigma},\qquad
 \dim_{K_0} V_\kappa\le m_\sigma.
\]
When the restricted share \(d_s\) is positive and \(T>s\),
\(\log d_s\le1041Z_b^3\). On the selected scale
\(P_d\le4164Z_b^4\).
\end{proposition}

\begin{proof}
The absolute value of each evaluation is bounded by the supremum. Conversely, by
\eqref{eq:slice-net-size}, \(\|F\|_{\sup,\sigma}\) is at most
four times the largest absolute value on the finite set, which is at most
\(\sqrt{m_\sigma}\,\|F\|_{2,\sigma}\). The finite evaluation
map is injective, so the dimension bound follows from the
number of complex coordinates.
The run's weight size \(|\kappa|\ge1\) makes this one bound
for the log of the squared-norm multiplier and for the log
of the dimension.

When \(d_s>0\) sections are selected by
restriction to the cut, the share bound is
\[
 \binom{T-1}{s}d_s\le R_j\le\dim_{K_0}V_\kappa.
\]
The selected range has \(T>s\), so
\(\log d_s\le1041Z_b^3\) before even using \(P_d\).
For \(d_s=0\) its determinant height contribution is zero.
This supplies the count in the height bound for the same
restricted short sections. It also supplies the metric and
rank bounds for the boundary determinant on the full cusp
lattice. With \(1+\log|\kappa|\le4Z_b\) on the selected
scale, the numerical absorption uses \(P_d\le4164Z_b^4\).

\end{proof}

\paragraph{Use of the evaluation norm in the index.}
The distinguished-prime construction selects the light,
short and auxiliary sections and the indicated restricted
basis using the running scalars. The bounds on the local
charts and lifts then apply to this family at all finite
places. Replacing the auxiliary Hermitian norm by the
finite evaluation norm carries these very sections, points,
trivializations and branch orders through the comparison.
For each embedding \(\sigma\) used for evaluation, the cost of using the finite evaluations in place of the supremum is controlled by
\[
  \|F\|_{2,\sigma}\le\|F\|_{\sup,\sigma}\le4\sqrt{m_\sigma}\,\|F\|_{2,\sigma}.
\]
The degrees of the projective cut and auxiliary section
still use their specified metrics, and the heights of the
generators still use the supremum. The finite norm enters the
boundary determinant by the inequality just proved. This is
the norm and count input to the four-term remainder below,
on each orientation of the run separately.
\label{cor:slice-evaluation-net-7}
